\documentclass[11pt]{article}
\usepackage[margin=1in]{geometry}
\usepackage{amsmath, amssymb, amsthm, mathtools}
\usepackage{booktabs}
\usepackage{enumitem}
\usepackage[colorlinks=true, linkcolor=blue, urlcolor=blue]{hyperref}
\usepackage{array}
\usepackage{microtype}

\setlength{\parskip}{0.6em}
\setlength{\parindent}{0pt}

\newcommand{\E}{\mathbb{E}}
\newcommand{\Var}{\mathrm{Var}}
\newcommand{\Cov}{\mathrm{Cov}}
\newcommand{\Prob}{\mathbb{P}}
\newcommand{\1}{\mathbf{1}}
\newcommand{\Bern}{\mathrm{Bern}}
\newcommand{\Bin}{\mathrm{Bin}}
\newcommand{\Geom}{\mathrm{Geom}}
\newcommand{\NBin}{\mathrm{NBin}}
\newcommand{\Pois}{\mathrm{Pois}}
\newcommand{\Unif}{\mathrm{Unif}}
\newcommand{\Normal}{\mathcal{N}}
\newcommand{\floor}[1]{\left\lfloor #1 \right\rfloor}
\newcommand{\ceil}[1]{\left\lceil #1 \right\rceil}

\title{Stat 110 Midterm Study Guide}
\author{David Wang}
\date{March 2026}

\begin{document}
\maketitle
\tableofcontents

\section{Scope and parameter conventions}
This guide is designed for the \emph{first half} of Stat 110, through the midterm review. I am using the following parameterizations throughout:
\begin{itemize}[leftmargin=1.5em]
    \item $X \sim \Geom(p)$ means $X$ is the \textbf{trial number of the first success}, so its support is $\{1,2,3,\dots\}$.
    \item $X \sim \NBin(r,p)$ means $X$ is the \textbf{trial number of the $r$th success}, so its support is $\{r,r+1,r+2,\dots\}$.
    \item $X \sim \Normal(\mu,\sigma^2)$ means mean $\mu$ and variance $\sigma^2$.
\end{itemize}

\section{Core definitions and identities}

\subsection{PMF, PDF, and CDF}
For a discrete random variable $X$, the \textbf{PMF} is
\[
    p_X(x) = \Prob(X=x).
\]
For a continuous random variable $X$, the \textbf{PDF} $f_X(x)$ satisfies
\[
    \Prob(a \le X \le b) = \int_a^b f_X(x)\,dx,
\]
and for all $x$,
\[
    f_X(x) = F_X'(x) \quad \text{when } F_X \text{ is differentiable at } x.
\]
For any random variable, the \textbf{CDF} is
\[
    F_X(x) = \Prob(X \le x).
\]
Useful facts:
\begin{itemize}[leftmargin=1.5em]
    \item $F_X(x)$ is non-decreasing.
    \item $\lim_{x\to -\infty}F_X(x)=0$ and $\lim_{x\to \infty}F_X(x)=1$.
    \item For discrete $X$, the CDF is a step function.
\end{itemize}

\subsection{Expectation and variance}
For discrete $X$,
\[
    \E[X] = \sum_x x\,\Prob(X=x).
\]
For continuous $X$,
\[
    \E[X] = \int_{-\infty}^{\infty} x f_X(x)\,dx.
\]
Variance is
\[
    \Var(X)=\E[(X-\E[X])^2]=\E[X^2]-(\E[X])^2.
\]
Standard deviation is $\sqrt{\Var(X)}$.

\subsection{LOTUS}
The \textbf{Law of the Unconscious Statistician (LOTUS)} lets you compute $\E[g(X)]$ \textit{without} first finding the distribution of $g(X)$. You do so using the PMF or PDF of random variable $X$.

For discrete $X$,
\[
    \E[g(X)] = \sum_x g(x)\,\Prob(X=x).
\]
For continuous $X$,
\[
    \E[g(X)] = \int_{-\infty}^{\infty} g(x)f_X(x)\,dx.
\]
This is useful for calculating variance:
\[
    \E[X^2] = \sum_x x^2\Prob(X=x) \quad \text{or} \quad \int x^2 f_X(x)\,dx.
\]

\subsection{Linearity of expectation}
For any random variables $X_1,\dots,X_n$ and constants $a_1,\dots,a_n$,
\[
    \E\left[\sum_{i=1}^n a_i X_i\right] = \sum_{i=1}^n a_i\E[X_i].
\]
\textbf{Important:} independence is \textit{not} required.

In particular,
\[
    \E[aX+b] = a\E[X]+b.
\]

\subsection{Variance rules}
\[
    \Var(aX+b)=a^2\Var(X).
\]
If $X$ and $Y$ are independent, then
\[
    \Var(X+Y)=\Var(X)+\Var(Y).
\]
Without independence, \textit{do not} separate the variance like that.

\subsection{Indicator random variables}
For an event $A$, let $I_A = 1$ when the event occurs and $0$ otherwise. It follows that 
\[
    \E[I_A]=\Prob(A).
\]
This is one of the most useful tricks. 

\subsection{Conditioning basics worth keeping on the sheet}
Conditional probability:
\[
    \Prob(A\mid B)=\frac{\Prob(A\cap B)}{\Prob(B)} \qquad (\Prob(B)>0).
\]

\subsection{Law of total probability}
For disjoint events $B_1,\dots,B_n$ forming a partition:
\[
    \Prob(A)=\sum_{i=1}^n \Prob(A\mid B_i)\Prob(B_i).
\]

\subsection{Bayes' rule}
\[
    \Prob(B_j\mid A)=\frac{\Prob(A\mid B_j)\Prob(B_j)}{\sum_{i=1}^n \Prob(A\mid B_i)\Prob(B_i)}.
\]

\subsection{Universality of the Uniform}
If $X$ has a continuous CDF $F$, then
\[
    F(X) \sim \Unif(0,1).
\]
Conversely, if $U\sim \Unif(0,1)$ and $F$ is a continuous strictly increasing CDF, then
\[
    F^{-1}(U)
\]
has CDF $F$.

\subsection{Location and scale}
If $Y=aX+b$, then
\[
    \E[Y]=a\E[X]+b, \qquad \Var(Y)=a^2\Var(X).
\]
For a Normal random variable,
\[
    X\sim \Normal(\mu,\sigma^2)
    \quad \Longrightarrow \quad
    \frac{X-\mu}{\sigma} \sim \Normal(0,1).
\]

\section{Distributions studied so far}

\subsection{Bernoulli distribution}
Let $X\sim \Bern(p)$ with support $\{0,1\}$. A Bernoulli trial has a result $1$ and probability $p$ of success.

\textbf{PMF}
\[
    \Prob(X=1)=p, \qquad \Prob(X=0)=1-p.
\]

\textbf{CDF}
\[
    F_X(x)=
    \begin{cases}
        0, & x<0,\\
        1-p, & 0\le x<1,\\
        1, & x\ge 1.
    \end{cases}
\]

\textbf{Mean and variance}
\[
    \E[X]=p, \qquad \Var(X)=p(1-p).
\]

\subsection{Binomial distribution}
Let $X\sim \Bin(n,p)$, the number of successes in $n$ independent Bern$(p)$ trials.

\textbf{Support}
\[
    x\in\{0,1,2,\dots,n\}.
\]

\textbf{PMF}
\[
    \Prob(X=x)=\binom{n}{x}p^x(1-p)^{n-x}.
\]

\textbf{CDF}
\[
    F_X(x)=\sum_{k=0}^{\floor{x}} \binom{n}{k}p^k(1-p)^{n-k}.
\]
(Interpret the sum as $0$ for $x<0$ and $1$ for $x\ge n$.)

\textbf{Mean and variance}
\[
    \E[X]=np, \qquad \Var(X)=np(1-p).
\]

\subsection{Hypergeometric distribution}
Let $X\sim \mathrm{HGeom}(N,K,n)$, where a population has size $N$, with $K$ successes, and we draw $n$ items \emph{without replacement}. Then $X$ is the number of successes drawn.

\textbf{Support}
\[
    x\in \left\{\max(0,n-(N-K)),\dots,\min(n,K)\right\}.
\]

\textbf{PMF}
\[
    \Prob(X=x)=\frac{\binom{K}{x}\binom{N-K}{n-x}}{\binom{N}{n}}.
\]

\textbf{CDF}
\[
    F_X(x)=\sum_{k\le \floor{x}} \frac{\binom{K}{k}\binom{N-K}{n-k}}{\binom{N}{n}},
\]
where the sum runs over allowable support values.

\textbf{Mean and variance}
\[
    \E[X]=n\frac{K}{N},
\]
\[
    \Var(X)=n\frac{K}{N}\left(1-\frac{K}{N}\right)\frac{N-n}{N-1}.
\]

\subsection{Discrete uniform distribution}
A common version is $X\sim \Unif\{1,2,\dots,n\}$, meaning all values are equally likely.

\textbf{PMF}
\[
    \Prob(X=x)=\frac{1}{n}, \qquad x\in\{1,2,\dots,n\}.
\]

\textbf{CDF}
\[
    F_X(x)=
    \begin{cases}
        0, & x<1,\\
        \dfrac{\floor{x}}{n}, & 1\le x<n,\\
        1, & x\ge n.
    \end{cases}
\]

\textbf{Mean and variance}
\[
    \E[X]=\frac{n+1}{2}, \qquad \Var(X)=\frac{n^2-1}{12}.
\]

\subsection{Geometric distribution}
Let $X\sim \Geom(p)$ be the trial number of the first success.

\textbf{Support}
\[
    x\in\{1,2,3,\dots\}.
\]

\textbf{PMF}
\[
    \Prob(X=x)=(1-p)^{x-1}p.
\]

\textbf{CDF}
\[
    F_X(x)=
    \begin{cases}
        0, & x<1,\\
        1-(1-p)^{\floor{x}}, & x\ge 1.
    \end{cases}
\]

\textbf{Mean and variance}
\[
    \E[X]=\frac{1}{p}, \qquad \Var(X)=\frac{1-p}{p^2}.
\]

\textbf{Special fact: memoryless property}
\[
    \Prob(X>m+n\mid X>m)=\Prob(X>n).
\]

\subsection{Negative binomial distribution}
Let $X\sim \NBin(r,p)$ be the trial number of the $r$th success.

\textbf{Support}
\[
    x\in\{r,r+1,r+2,\dots\}.
\]

\textbf{PMF}
\[
    \Prob(X=x)=\binom{x-1}{r-1}p^r(1-p)^{x-r}.
\]

\textbf{CDF}
\[
    F_X(x)=\sum_{k=r}^{\floor{x}} \binom{k-1}{r-1}p^r(1-p)^{k-r}.
\]

\textbf{Mean and variance}
\[
    \E[X]=\frac{r}{p}, \qquad \Var(X)=\frac{r(1-p)}{p^2}.
\]

\textbf{Connection}
When $r=1$, the negative binomial becomes the geometric distribution.

\subsection{Poisson distribution}
Let $X\sim \Pois(\lambda)$.

\textbf{Support}
\[
    x\in\{0,1,2,3,\dots\}.
\]

\textbf{PMF}
\[
    \Prob(X=x)=e^{-\lambda}\frac{\lambda^x}{x!}.
\]

\textbf{CDF}
\[
    F_X(x)=\sum_{k=0}^{\floor{x}} e^{-\lambda}\frac{\lambda^k}{k!}.
\]

\textbf{Mean and variance}
\[
    \E[X]=\lambda, \qquad \Var(X)=\lambda.
\]

\textbf{Approximation idea}

A Binomial$(n,p)$ random variable is often well approximated by Poisson$(\lambda=np)$ when $n$ is large and $p$ is small.

\subsection{Continuous uniform distribution}
Let $X\sim \Unif(a,b)$ with $a<b$.

\textbf{PDF}
\[
    f_X(x)=
    \begin{cases}
        \dfrac{1}{b-a}, & a\le x\le b,\\
        0, & \text{otherwise}.
    \end{cases}
\]

\textbf{CDF}
\[
    F_X(x)=
    \begin{cases}
        0, & x<a,\\
        \dfrac{x-a}{b-a}, & a\le x\le b,\\
        1, & x>b.
    \end{cases}
\]

\textbf{Mean and variance}
\[
    \E[X]=\frac{a+b}{2}, \qquad \Var(X)=\frac{(b-a)^2}{12}.
\]

\subsection{Normal distribution}
Let $X\sim \Normal(\mu,\sigma^2)$ with $\sigma>0$.

\textbf{PDF}
\[
    f_X(x)=\frac{1}{\sigma\sqrt{2\pi}}e^\frac{-(x-\mu)^2}{2\sigma^2}.
\]

\textbf{CDF}
There is no elementary closed form. It is written as
\[
    F_X(x)=\Phi\!\left(\frac{x-\mu}{\sigma}\right),
\]
where $\Phi$ is the standard Normal CDF.

\textbf{Mean and variance}
\[
    \E[X]=\mu, \qquad \Var(X)=\sigma^2.
\]

\textbf{Standardization}
\[
    Z=\frac{X-\mu}{\sigma} \sim \Normal(0,1).
\]

\section{A few especially important problem-solving tools}

\subsection{CDF method}
For any random variable,
\[
    \Prob(a<X\le b)=F_X(b)-F_X(a)
\]
with the endpoint details adjusted appropriately for discrete variables.

\subsection{Turning a hard variable into indicators}
If $X$ counts how many of several events occur, write
\[
    X=I_1+I_2+\cdots+I_n.
\]
Then
\[
    \E[X]=\E[I_1]+\cdots+\E[I_n]=\Prob(A_1)+\cdots+\Prob(A_n).
\]
This is central for occupancy, matching, and coupon-collector-type arguments.

\subsection{Symmetry}
When a problem is symmetric across labels, positions, or people, look for a way to avoid direct counting and instead argue that certain probabilities or expectations must match by symmetry.

\subsection{When to condition}
If a direct computation looks messy, try conditioning on a strategically chosen event or random variable that breaks the problem into simpler pieces.

\section{Handy formulas}
\begin{align*}
\E[aX+b] &= a\E[X]+b,\\
\Var(X) &= \E[X^2]-(\E[X])^2,\\
\Var(aX+b) &= a^2\Var(X),\\
\E[I_A] &= \Prob(A),\\
\E[g(X)] &= \sum_x g(x)\Prob(X=x) \quad \text{or} \quad \int g(x)f_X(x)\,dx,\\
F_X(x) &= \Prob(X\le x),\\
X\sim \Normal(\mu,\sigma^2) &\Rightarrow \frac{X-\mu}{\sigma}\sim \Normal(0,1),\\
X\sim \Geom(p) &\Rightarrow \E[X]=\frac{1}{p},\\
X\sim \Pois(\lambda) &\Rightarrow \E[X]=\Var(X)=\lambda.
\end{align*}

\end{document}
